\begin{center}\emph{Forms with negative determinant.}\end{center}

\Article{171}

\textsc{Problem.}\quad \emph{Given any form $(a,b,a')$ with negative determinant $=-D$, where $D$ denotes a positive number, find a properly equivalent form $(A,B,C)$ in which $A$ is no greater than $\surd\tfrac43D,$ $C$, and no smaller than $2B$.}

\emph{Solution.}\quad We assume that the given form does not satisfy all three conditions at once; otherwise there would be no need to look for another form.\quad Let $b'$ be the absolutely least residue of the number $-b$ modulo $a'$\HistoricalNote{*)}{Notice that if, in a form $(a,b,a')$, the first or last term $a$ or $a'$ is $=0$, its determinant is a positive square, so this cannot happen in the present case. --- For a similar reason, the outer terms $a,a'$ of a form with negative determinant cannot have opposite signs.}, and let $a''=\frac{b'b'+D}{a'}$, which is an integer because $b'b'\equiv bb$, $b'b'+D\equiv bb+D\equiv aa'\equiv$ $0(\operatorname{mod.}a').$\quad Now if $a''<a'$, again let $b''$ be the absolutely least residue of $-b'$ modulo $a''$, and let $a'''=\frac{b''b''+D}{a''}$. If again $a'''<a''$, once more let $b'''$ be the absolutely least residue of $-b''$ modulo $a'''$ and let $a''''=\frac{b'''b'''+D}{a'''}$.\quad Continue this operation until the sequence $a',a''$, $a''',a''''$ etc. reaches a term $a^{m+1}$ no smaller than its predecessor $a^m$. This must eventually happen, since otherwise we would have an infinite sequence of integers that keep decreasing.\quad Then the form $(a^m,b^m,a^{m+1})$ satisfies all the conditions.

\emph{Proof.}\quad I.\quad In the sequence of forms $(a,b,a'),\ (a',b',a''),\ (a'',b'',a''')$ etc., each is a neighbor of its predecessor, so the last is properly equivalent to the first (Articles $159$, $160).$
\SourceLeaf{152}
II. Since $b^m$ is the absolutely least residue of $-b^{m-1}$ modulo $a^m$, it is no greater than $\frac12 a^m$ (Article 4).

III. Since $a^m a^{m+1}=D+b^m b^m$, and $a^{m+1}$ is not $<a^m$, $a^m a^m$ is not $>D+b^m b^m$. Since $b^m$ is not $>\frac12 a^m$, $a^m a^m$ is not $>D+\frac14 a^m a^m$, and $\frac34 a^m a^m$ is not $>D$; finally, $a^m$ is \emph{not} $>\surd{\frac43 D}$.

\emph{Example.} Let the given form be $(304,217,155)$, whose determinant $=-31$. Here we find the sequence of forms:
\[
(304,217,155),\ (155,-62,25),\ (25,12,7),\ (7,2,5),\ (5,-2,7).
\]
The last is the one we want. --- In the same way, for the form $(121,49,20)$, whose determinant $=-19$, we find the equivalent forms $(20,-9,5)$, $(5,-1,4)$, $(4,1,5)$; so $(4,1,5)$ is the form we want.

We will call forms $(A,B,C)$ with negative determinant, in which $A$ is no greater than $\surd{\frac43 D},C$ and no smaller than $2B$, \emph{reduced forms}. So for any form with negative determinant, a properly equivalent reduced form can be found.

\Needspace{8\baselineskip}
\Article{172}
\textsc{Problem.} \emph{Find the conditions under which two nonidentical reduced forms with the same determinant $-D$, $(a,b,c)$, $(a',b',c')$, can be properly equivalent.}

\emph{Solution.} We may assume that $a'$ is not $>a$, and that the form $axx+2bxy+cyy$ changes into $a'x'x'+2b'x'y'+c'y'y'$ under the proper substitution $x=\alpha x'+\beta y'$, $y=\gamma x'+\delta y'$. Then we have the equations
\[
\begin{aligned}
a\alpha\alpha+2b\alpha\gamma+c\gamma\gamma&=a'&\quad\mathinner{.\,.\,.\,.\,.\,.\,.\,.}\quad&[1]\\
a\alpha\beta+b(\alpha\delta+\beta\gamma)+c\gamma\delta&=b'&\quad\mathinner{.\,.\,.\,.\,.\,.\,.}\quad&[2]\\
\alpha\delta-\beta\gamma&=1&\quad\mathinner{.\,.\,.\,.\,.\,.\,.\,.\,.\,.\,.}\quad&[3]
\end{aligned}
\]
From [1] we get $aa'=(a\alpha+b\gamma)^2+D\gamma\gamma$, so $aa'$ is positive. Since $ac=D+bb$, $a'c'=D+b'b'$, $ac$ and $a'c'$ are also positive; so $a,a',c,c'$ all have the same sign. But both $a$ and $a'$ are not $>\surd{\frac43 D}$, so $aa'$ is not $>\surd{\frac43 D}$; still less, then, can $D\gamma\gamma\ (=aa'-(a\alpha+b\gamma)^2)$ be greater than $\frac43 D$. Thus $\gamma$ is either $=0$ or $=\pm1$.

I. If $\gamma=0$, [3] gives either $\alpha=1.\,\delta=1$ or $\alpha=-1$, $\delta=-1$.
\SourceLeaf{153}
In both cases [1] gives $a'=a$, and [2] gives $b'-b=\pm\beta a$. But $b$ is not $>\frac12 a$, and $b'$ is not $>\frac12 a'$, so it is also not $>\frac12 a$. Thus the equation $b'-b=\pm\beta a$ can hold only if

\emph{either} $b=b'$, which would give $c'=\frac{b'b'+D}{a'}=\frac{bb+D}{a}=c$, making the forms $(a,b,c)$, $(a',b',c')$ identical, against our assumption;

\emph{or} $b=-b'=\pm\frac12 a$. In this case too, $c'=c$, and the form $(a',b',c')$ is $(a,-b,c)$, \emph{that is}, the opposite of the form $(a,b,c)$. It is also clear that these forms are ambiguous, since $2b=\pm a$.

II. If $\gamma=\pm1$, [1] gives $a\alpha\alpha+c-a'=\pm2b\alpha$. But $c$ is no smaller than $a$, so it is no smaller than $a'$. Thus $a\alpha\alpha+c-a'$, or $2b\alpha$, is certainly no smaller than $a\alpha\alpha$. Since $2b$ is no greater than $a$, $\alpha$ is no smaller than $\alpha\alpha$; so necessarily either $\alpha=0$ or $=\pm1$.

1) If $\alpha=0$, [1] gives $a'=c$, and since $a$ is neither greater than $c$ nor smaller than $a'$, necessarily $a'=a=c$. Further, [3] gives $\beta\gamma=-1$, so [2] gives $b+b'=\pm\delta c=\pm\delta a$. From this, as in (I), we get

\emph{either} $b=b'$, in which case the forms $(a,b,c)$, $(a',b',c')$ would be identical, against our assumption;

\emph{or} $b=-b'$, in which case the forms $(a,b,c)$, $(a',b',c')$ are opposite.

2) If $\alpha=\pm1$, [1] gives $\pm2b=a+c-a'$. So, since $a$ is not smaller and $c<a'$ does not hold, $2b$ is not $<a$ and not $<c$. But $2b$ is also not $>a$ or $>c$, so necessarily $\pm2b=a=c$, and from the equation $\pm2b=a+c-a'$, also $=a'$. Thus [2] gives
\[
b'=a(\alpha\beta+\gamma\delta)+b(\alpha\delta+\beta\gamma)
\]
or, since $\alpha\delta-\beta\gamma=1$,
\[
b'-b=a(\alpha\beta+\gamma\delta)+2b\beta\gamma=a(\alpha\beta+\gamma\delta\pm\beta\gamma)
\]
so necessarily, as before,

\emph{either} $b=b'$, making the forms $(a,b,c)$, $(a',b',c')$ identical, against our assumption;

\emph{or} $b=-b'$, so those forms are opposite. In this case too, since $a=\pm2b$, the forms are ambiguous.

From all this we conclude that the forms $(a,b,c)$, $(a',b',c')$ can be properly equivalent only if they are opposite and, at the same time, \emph{either} ambiguous \emph{or} $a=c=a'=c'$.
\SourceLeaf{154}
That the forms $(a,b,c)$, $(a',b',c')$ are properly equivalent in these cases could also have been seen easily in advance. For opposite forms must be improperly equivalent, and if they are also ambiguous, they must be properly equivalent too. If $a=c$, the form $(\frac{D+(a-b)^2}{a},a-b,a)$ is a neighbor of the form $(a,b,c)$ and so equivalent to it. But since $D+bb=ac=aa$, we have $\frac{D+(a-b)^2}{a}=2a-2b$, and the form $(2a-2b,a-b,a)$ is ambiguous; so $(a,b,c)$ is properly equivalent to its opposite too.

We can now decide just as easily when two reduced forms $(a,b,c)$, $(a',b',c')$ that are not opposite can be improperly equivalent. They are improperly equivalent if $(a,b,c)$, $(a',-b',c')$, which are not identical, are properly equivalent, and conversely. Thus the condition for \emph{the former} to be improperly equivalent is that they are identical and, in addition, either ambiguous or $a=c$. --- Reduced forms that are neither identical nor opposite can be neither properly nor improperly equivalent.

\Needspace{8\baselineskip}
\Article{173}
\textsc{Problem.} \emph{Given two forms $F$ and $F'$ with the same negative determinant, determine whether they are equivalent.}

\emph{Solution.} Find two reduced forms $f$, $f'$ properly equivalent to the forms $F$, $F'$, respectively. If the forms $f$, $f'$ are properly equivalent, improperly equivalent, or equivalent in both ways, so are $F$, $F'$. If $f$, $f'$ are not equivalent in any way, neither are $F$, $F'$.

The preceding article gives four possible cases:

1) If $f$, $f'$ are neither identical nor opposite, $F$, $F'$ are not equivalent in any way.

2) If $f$, $f'$ are \emph{first} either identical or opposite, and \emph{second} either ambiguous or have equal outer terms, $F$, $F'$ are both properly and improperly equivalent.

3) If $f$, $f'$ are identical, but neither ambiguous nor have equal outer terms, $F,F'$ are only properly equivalent.

4) If $f$, $f'$ are opposite, but neither ambiguous nor have equal outer terms, $F,F'$ are only improperly equivalent.

\emph{Example.} For the forms $(41,35,30)$, $(7,18,47)$, whose determinant $=-5$, we find the equivalent reduced forms $(1,0,5)$, $(2,1,3)$; so the given forms are not equivalent in any way. --- But the forms $(23,38,63)$, $(15,20,27)$ have the same equivalent reduced
\SourceLeaf{155}
form $(2,1,3)$. Since it is also ambiguous, the forms $(23,38,63)$, $(15,20,27)$ are both properly and improperly equivalent. --- For the forms $(37,53,78)$, $(53,73,102)$, the equivalent reduced forms are $(9,2,9)$, $(9,-2,9)$. Since the reduced forms are opposite and their outer terms are equal, the given forms are both properly and improperly equivalent.

\Needspace{8\baselineskip}
\Article{174}
The number of all reduced forms with a given determinant $-D$ is always finite and fairly small in relation to the number $D$. The forms themselves can be found in two ways. Denote the reduced forms with determinant $-D$ generally by $(a,b,c)$; we then have to determine all the values of $a,b,c$.

\emph{First method.} For $a$, take all numbers, both positive and negative, no greater than $\surd{\frac43 D}$, for which $-D$ is a quadratic residue. For each $a$, let $b$ successively equal all values of the expression $\surd{-D}\ (\mathrm{mod.}\ a)$ no greater than $\frac12 a$, taken both positively and negatively. For $c$, given each determined pair of values of $a,b$, put $=\frac{D+bb}{a}$. Discard any forms produced this way in which $c<a$; the others are clearly reduced.

\emph{Second method.} For $b$, take all numbers, both positive and negative, no greater than $\frac12\surd{\frac43 D}$, or $\surd{\frac13 D}$. For each $b$, factor $bb+D$ in every possible way into two factors (also taking differences of sign into account), both no smaller than $2b$. Put one factor---the smaller one when they are unequal---$=a$, and the other $=c$. Since $a$ is not $>\surd{\frac43 D}$, all forms obtained this way are clearly reduced. --- Finally, it is clear that no reduced form can fail to be found by either method.

\emph{Example.} Let $D=85$. Here the bound on the values of $a$ is $\surd{\frac{340}{3}}$, which lies between $10$ and $11$. The numbers between $1$ and $10$ (inclusive) for which $-85$ is a residue are $1,2,5,10$. This gives twelve forms:
\[
\begin{gathered}
(1,0,85),\ (2,1,43),\ (2,-1,43),\ (5,0,17),\\
(10,5,11),\ (10.\,-5,11);\ (-1,0,-85),\ (-2,1,-43),\\
(-2,-1,-43).\ (-5,0,-17),\ (-10,5,-11,\ (-10,-5,-11).
\end{gathered}
\]
The other method gives, for the values of $b$, the bound $\surd{\frac{85}{3}}$, which lies between $5$ and $6$. For $b=0$, we get the forms
\SourceLeaf{156}
\[
(1,0,85),\ (-1,0,-85),\ (5,0,17),\ (-5,0,-17),
\]
and for $b=\pm1$, these: $(2,\pm1,43)$, $(-2,\pm1,-43)$.

For $b=\pm2$, there are none, since $89$ cannot be factored into two factors both not $<4$. The same holds for $\pm3$, $\pm4$. Finally, for $b=\pm5$, we get
\[
(10,\pm5,11),\ (-10,\pm5.\,-11).
\]

\Needspace{8\baselineskip}
\Article{175}
From all reduced forms with a given determinant, discard one of each pair of forms that, though not identical, are properly equivalent. The remaining forms have this notable property: every form with the same determinant is properly equivalent to one of them, and to only one (otherwise some of them would be properly equivalent to one another). Thus \emph{all forms with the same determinant can be divided into as many classes as there are forms remaining}, by putting forms properly equivalent to the same reduced form into the same class. For $D=85$, for example, the remaining forms are
\[
\begin{gathered}
(1,0,85),\ (2,1,43),\ (5,0,17),\ (10,5,11),\\
(-1,0.\,-84),\ (-2,1,-43).\ (-5,0,-17),\ (-10,5,-11)
\end{gathered}
\]
so all forms with determinant $-85$ can be divided into eight classes, according as they are properly equivalent to the first form, the second, etc. Clearly, forms in the same class are properly equivalent, and forms from different classes cannot be properly equivalent. We will discuss this subject of classifying forms much more fully below. Here we add just one observation. We showed above that if the determinant of the form $(a,b,c)$ is negative $=-D$, $a$ and $c$ have the same sign (since $ac=bb+D$, which is positive). In the same way, it is easy to see that if the forms $(a,b,c)$, $(a',b',c')$ are equivalent, all of $a,c,a',c'$ have the same sign. For if the first changes into the second under the substitution $x=\alpha x'+\beta y'$, $y=\gamma x'+\delta y'$, then $a\alpha\alpha+2b\alpha\gamma+c\gamma\gamma=a'$, giving $aa'=(a\alpha+b\beta)^2+D\gamma\gamma$, which is certainly nonnegative. Since $a$ cannot be zero and $a'=0$ cannot hold, $aa'$ is positive, so $a,a'$ have the same sign.

This makes it clear that forms whose outer terms are positive are completely separate from those whose outer terms are negative. It is enough to consider just the reduced forms with positive outer terms, since
\SourceLeaf{157}
the others are equal in number and come from them by giving the outer terms opposite signs. The same holds for the forms to be discarded from the reduced forms and those to be retained.

\Needspace{8\baselineskip}
\Article{176}
Here, then, is a table of forms for certain negative determinants, according to which all other forms with the same determinant can be divided into classes. In line with the observation in the preceding article, we give only half of them, those with positive outer terms.
\[
\begin{array}{r|l}
\multicolumn{1}{r}{D}&\\
1&(1,0,1).\\
2&(1,0,2).\\
3&(1,0,3),\ (2,1,2).\\
4&(1,0,4),\ (2,0,2).\\
5&(1,0,5),\ (2,1,3).\\
6&(1,0,6),\ (2,0,3).\\
7&(1,0,7),\ (2,1,4).\\
8&(1,0,8),\ (2,0,4),\ (3,1,3).\\
9&(1,0,9),\ (2,1,5),\ (3.\,0,3).\\
10&(1,0,10),\ (2,0,5).\\
11&(1,0,11),\ (2,1,6),\ (3,1,4),\ (3,-1,4).\\
12&(1,0,12),\ (2,0,6),\ (3,0,4),\ (4,2,4).
\end{array}
\]
There would be no need to continue this table further here, since below we will show how to arrange it much more usefully.

So every form with determinant $-1$ is properly equivalent to the form $xx+yy$ if its outer terms are positive, or to $-xx-yy$ if they are negative; every form with determinant $-2$ whose outer terms are positive is properly equivalent to the form $xx+2yy$, etc.; every form with determinant $-11$ whose outer terms are positive, to one of $xx+11yy$, $2xx+2xy+6yy$, $3xx+2xy+4yy$, $3xx-2xy+4yy$, etc.
